% !PRINTED 479
\weberpage{479}

Hierauf können wir, wenn $\varrho$ durch die Wurzeln $\alpha$ ausgedrückt ist, nach Voraussetzung die Permutation $\pi$ anwenden und erhalten also, wenn dadurch $\varrho$ in $\varrho_a$ übergeht, $g(\varrho_a) = 0$; d.\,h. $\varrho_a$ ist auch eine Wurzel der Resolvente, und die Permutation $\pi$ entspricht also einer Substitution $(\varrho, \varrho_a)$ des Körpers $\Omega(\varrho)$, d.\,h. $\pi$ gehört zur \gesperrt{Galois'schen Gruppe}.

Daraus schliessen wir noch auf folgenden wichtigen Satz:

d) \gesperrt{Ist $P$ eine Gruppe von Permutationen der $m$ Wurzeln $\alpha$, der die Eigenschaften a) und b) zukommen, so ist $P$ die Galois'sche Gruppe der Gleichung $F(x) = 0$.}

Denn zunächst gehört nach c) jede Permutation von $P$ der \gesperrt{Galois'schen Gruppe} an, und $P$ ist also gewiss ein Theiler von dieser.

Wenn aber $P$ nur $\nu$ Permutationen umfasst, so mögen diese mit
\weq{5}{\pi_1, \pi_2 \ldots \pi_\nu}
bezeichnet sein. Wenden wir diese Permutationen auf $\varrho$ an, so möge sich ergeben:
\weq{6}{\varrho_1, \varrho_2 \ldots \varrho_\nu.}
Wenden wir eine der Permutationen (5), etwa $\pi_k$, auf eine der Grössen (6), etwa auf $\varrho_i$ an, so ist der Erfolg derselbe, als ob $\pi_i \pi_k$ auf $\varrho$ angewandt worden sei; das Ergebniss dieser Permutation soll $\varrho_i'$ sein. Nun liegt aber in der Voraussetzung, dass $P$ eine Gruppe sei, dass auch die componirte Permutation $\pi_i \pi_k$ zu $P$ gehört, dass also $\varrho_i'$ unter den Grössen (6) enthalten sei. Zugleich ist $\varrho_i'$ nach \S.\,146, 2. von $\varrho_h'$ verschieden, sobald $\varrho_i$ von $\varrho_h$ verschieden ist. Daraus folgt, dass die Grössen
\weq{7}{\varrho_1', \varrho_2' \ldots \varrho_\nu'}
mit den Grössen (6), von der Ordnung abgesehen, übereinstimmen.

Das Product
\[
g'(t) = (t - \varrho_1)(t - \varrho_2) \ldots (t - \varrho_\nu)
\]
bleibt also durch die Permutationen der Gruppe $P$ ungeändert, und ist folglich, da wir die Eigenschaft b) von $P$ voraussetzen, eine rationale Function von $t$ in $\Omega$. Zugleich ist $g'(t)$ ein Theiler von $g(t)$ und muss daher, da $g(t)$ irreducibel ist, mit $g(t)$ übereinstimmen, also ist $\nu$ nicht kleiner als der Grad der \gesperrt{Galois'schen Gruppe} und $P$ ist mit der \gesperrt{Galois'schen Gruppe} identisch.
